\section{Legislação}\label{sec:legislacao}

\subsection{Diploma em vigor}

O \textbf{Decreto-Lei n.º 54/2005, de 3 de março}\footnote{\url{https://files.diariodarepublica.pt/1s/2005/03/044a00/18921897.pdf}}
aprova o \emph{Regulamento do Número e Chapa de Matrícula dos Automóveis, Seus
Reboques, Motociclos, Triciclos e Quadriciclos de Cilindrada Superior a
50~cm³}. Os artigos principais do Regulamento são:

\begin{description}
  \item[Art. 3.º, número de matrícula.] Dois grupos de dois algarismos e um grupo
    de duas letras, separados por traços. A posição das letras depende da data:
    \plate{AA-00-00} até 29/02/1992, \plate{00-00-AA} a partir de 01/03/1992 e
    \plate{00-AA-00} depois de esgotada a série anterior. O n.º~3 já previa que,
    esgotada a série \plate{00-AA-00}, o número passaria a ter dois grupos de duas
    letras e um grupo central de dois algarismos (\plate{AA-00-AA}).
  \item[Art. 4.º, reboques e exportação.] Os reboques têm uma ou duas letras do
    serviço regional (anexo~I) seguidas de um número de ordem; os veículos para
    exportação, um número de ordem seguido da inicial de Lisboa, Porto, Açores ou
    Madeira.
  \item[Art. 5.º, chapa de matrícula.] O modelo depende da data da matrícula:
    fundo preto e caracteres brancos até 31/12/1991 (anexo~II); retrorrefletora,
    de fundo branco com caracteres e rebordo pretos, de 01/01/1992 a 31/12/1997
    (anexo~III); igual, com o ano e o mês da primeira matrícula na extremidade
    direita, a partir de 01/01/1998 (anexo~IV). As chapas antigas podem ser
    substituídas por modelos posteriores.
  \item[Art. 6.º, casos particulares.] Exportação: chapa amarela com caracteres
    pretos (anexo~V). Corpo diplomático e consular, pessoal de missões
    estrangeiras e entidades abrangidas pelo Protocolo relativo aos Privilégios e
    Imunidades das Comunidades Europeias: caracteres, traços e rebordo vermelhos.
  \item[Art. 7.º, instalação.] Os automóveis têm duas chapas, uma à frente e
    outra atrás.
\end{description}

O próprio decreto-lei revoga os artigos 1.º a 3.º do Decreto Regulamentar
n.º~13/98 e os artigos 35.º e 37.º do Regulamento do Código da Estrada de 1954
(art.~5.º), e fixa em 01/01/2006 a entrada em vigor das regras sobre a
comercialização das chapas (art.~6.º).

\subsection{Alterações}

\begin{longtable}{@{}p{0.3\linewidth} p{0.64\linewidth}@{}}
\toprule
\textbf{Diploma} & \textbf{O que mudou} \\
\midrule\endhead
\bottomrule\endfoot
DL 106/2006, de 8 de junho & Alargou o regulamento aos ciclomotores e reviu as chapas dos motociclos. \\
DL 112/2009, de 18 de maio & 2.ª alteração; manteve as datas de 1991 e 1997 do art.~5.º. \\
Lei 46/2010, de 7 de setembro & 3.ª alteração. \\
DL 2/2020, de 14 de janeiro & 4.ª alteração. Pôs em uso o formato \plate{AA 00 AA}, retirou os traços e a barra do mês e do ano, e aproximou a chapa do modelo da maioria dos Estados-Membros da UE. Revogou os n.ºs~5 e~7 do art.~5.º do Regulamento. \\
DL 9/2021 & Consta da versão consolidada de 29/01/2021; o que mudou nas matrículas está por confirmar. \\
\end{longtable}

Segundo o IMT, o formato \plate{AA 00 AA} aplica-se só a veículos novos; o
dono de um veículo de uma série anterior pode pedir uma chapa sem o mês e o ano.
Não foi possível confirmar se houve alterações depois de 2021.

\subsection{Diplomas relacionados}

\begin{itemize}
  \item Decreto-Lei n.º 107/2006, de 8 de junho: regulamento de atribuição de
        matrícula das máquinas industriais (RAMMI).
  \item Despacho n.º 20301/2006: especificações técnicas das chapas (por confirmar).
  \item Lista de chapas de matrícula aprovadas pelo IMT (24/06/2026).
\end{itemize}

\subsection{Diplomas anteriores}

\begin{longtable}{@{}l p{0.38\linewidth} p{0.42\linewidth}@{}}
\toprule
\textbf{Ano} & \textbf{Diploma} & \textbf{O que estabeleceu} \\
\midrule\endhead
\bottomrule\endfoot
1901 & Decreto de 3 de outubro de 1901 & Primeira obrigação de uma chapa metálica atrás, com o número da licença e a sede do distrito. \\
1937 & Por identificar & Série nacional \plate{AA-00-00}, de fundo preto e caracteres brancos, desde 01/01/1937. \\
1954 & Regulamento do Código da Estrada (Decreto n.º 39\,987, de 22/12/1954) & Os arts.~35.º e~37.º regularam o número e a chapa de matrícula até 2005. \\
1991 & Portaria n.º 884/91, de 28 de agosto & Criou a chapa retrorrefletora de fundo branco e caracteres pretos. \\
1991 & Portaria n.º 1161/91 & Alterou o ponto~2 da Portaria 884/91 (efeitos desde 11/11/1991). \\
1992 & Aplicação da Portaria 884/91 & Passagem para \plate{00-00-AA} a 01/03/1992. \\
1997 & Portaria n.º 1174-A/97, de 17 de novembro & Atualizou as regras das chapas; terá introduzido a barra do ano e do mês usada desde 01/01/1998 (por confirmar). \\
1998 & Decreto Regulamentar n.º 13/98, de 15 de junho & Ciclomotores e motociclos até 50~cm³: três letras do município e número de ordem, fundo amarelo. Tratores agrícolas e seus reboques: fundo verde. \\
2005 & Decreto-Lei n.º 54/2005 & Juntou as regras num único regulamento. \\
\end{longtable}

Ficam por identificar o diploma que criou o sistema de 1937 e a divisão por
zonas, os diplomas dos sistemas de 1911 e 1918, e o conteúdo exato da Portaria
1174-A/97.
